\subsection{The ring Z((T))}

Let A be a ring. Endow the A-module \(\mathbf{A}^{\mathbf{Z}}\) with the product topology of the discrete topologies. The elements \((a_{n})\in\mathbf{A}^{\mathbf{Z}}\) such that there exists \(n_{0}\in\mathbf{Z}\) with \(a_{n}=0\) for \(n<n_{0}\) form a submodule B of \(\mathbf{A}^{\mathbf{Z}}\) . If \(a=(a_{n})\in\mathbf{B}\), \(b=(b_{n})\in\mathbf{B}\), and we set \(ab=c\) , with \(c_{n}=\sum_{i+j=n}a_{i}b_{j}\) , one defines on B a structure of A-algebra. Let T be the element \((\theta_{n})\) of B such that \(\theta_{n}=0\) for \(n\neq1\) and \(\theta_{1}=1\) . Then T is invertible in B; for every element \(a=(a_{n})\) of B, the family \((a_{n}\mathrm{T}^{n})_{n\in\mathbf{Z}}\) is summable in \(\mathbf{A}^{\mathbf{Z}}\) and one has

\[
a = \sum_ {n \in \mathbf {Z}} a _ {n} \mathrm{T} ^ {n}.
\]

In the rest of this chapter, we denote by A((T)) the A-algebra B; it contains as subalgebras the algebra A{[}{[}T{]}{]} of formal series and the algebra A{[}T, T \(^{-1}\) {]}; their intersection is the algebra A{[}T{]} of polynomials.

Remark. --- The ring A((T)) is naturally identified with the ring of fractions A{[}{[}T{]}{]} \(_{T}\) of the ring A{[}{[}T{]}{]} defined by the multiplicative subset consisting of the powers of T.

For \(n, p\) in \(\mathbf{Z}\) , we define the integer \(\begin{bmatrix} n \\ p \end{bmatrix}\) by

\[
\left\{ \begin{array}{l} {\left[ \begin{array}{l} n \\ p \end{array} \right] = 0 \quad \text {if} \quad p <   0 \quad \text {or} \quad p > n,} \\ {\left[ \begin{array}{l} n \\ p \end{array} \right] = \binom {n} {p} = \frac {n (n - 1) \dots (n - p + 1)}{p !} \quad \text {if} \quad 0 \leqslant p \leqslant n.} \end{array} \right.\tag{1}
\]

We have \(\begin{bmatrix} n \\ p \end{bmatrix} = \begin{bmatrix} n \\ n - p \end{bmatrix}\) for \(n, p \in \mathbf{Z}\) .

Lemma 1.--- The element 1 - T of \(\mathbf{Z}((\mathrm{T}))\) is invertible. For every integer \(r > 0\) we have

\[
(1 - \mathrm{T}) ^ {- r} = \sum_ {n \in \mathbf {Z}} \left[ \begin{array}{c} n + r - 1 \\ r - 1 \end{array} \right] \mathrm{T} ^ {n} = \sum_ {n \in \mathbf {N}} \binom {n + r - 1} {r - 1} \mathrm{T} ^ {n}.
\]

Indeed, 1 - T is invertible in the ring \(\mathbf{Z}[[\mathrm{T}]]\) , with inverse \(\sum_{m\geqslant 0}\mathbf{T}^m\) ; we thus have

\[
(1 - \mathrm{T}) ^ {- r} = (\sum_ {m \geqslant 0} \mathrm{T} ^ {m}) ^ {r} = \sum_ {m _ {1}, \dots , m _ {r} \geqslant 0} \mathrm{T} ^ {m _ {1} + m _ {2} + \dots + m _ {r}},
\]

and the announced formula follows from E, III, p. 44, prop. 15.

Let \(\mathrm{Q}(\mathrm{T}) \in \mathbf{Z}[\mathrm{T}, \mathrm{T}^{-1}]\) , \(r\) be an integer \(> 0\) , and \(\mathrm{F} = (1 - \mathrm{T})^{-r}\mathrm{Q} \in \mathbf{Z}((\mathrm{T}))\) . Put

\[
\mathrm{Q} (\mathrm{T}) = \sum_ {i \in \mathbf {Z}} a _ {i} \mathrm{T} ^ {i}, \quad \mathrm{F} = \sum_ {n \in \mathbf {Z}} \alpha_ {n} \mathrm{T} ^ {n}.
\]

Then, by Lemma 1, we have

\[
\alpha_ {n} = \sum_ {i \in \mathbf {Z}} a _ {i} \left[ \begin{array}{c} n - i + r - 1 \\ r - 1 \end{array} \right] = \sum_ {i \leqslant n} a _ {i} \binom {n - i + r - 1} {r - 1}.
\]

Let \(n_1\) be the supremum in \(\overline{\mathbf{R}}\) of the set of integers \(i \in \mathbf{Z}\) such that \(a_i \neq 0\) . For every integer \(n \geqslant n_1\) , one has \(\alpha_n = \tilde{\alpha}(n)\) , where \(\tilde{\alpha}\) is the polynomial of \(\mathbf{Q}[\mathbf{X}]\) defined by

\[
\tilde {\alpha} (X) = \frac {1}{(r - 1) !} \sum_ {i \in \mathbf {Z}} a _ {i} \prod_ {j = 1} ^ {r - 1} (X - i + j).
\]

If we set \(c = \mathbf{Q}(1) = \sum_{i \in \mathbf{Z}} a_i\) , one has \(\tilde{\alpha}(\mathbf{X}) = c\mathbf{X}^{r-1}/(r - 1)! + \theta(\mathbf{X})\) , where \(\theta\) is a polynomial of degree \(\leqslant r - 2\) . Consequently, we have

\[
\alpha_ {n} = c \frac {n ^ {r - 1}}{(r - 1) !} + \rho_ {n} n ^ {r - 2},\tag{2}
\]

where the rational number \(\rho_{n}\) tends to a limit as \(n\) increases indefinitely. We deduce the relation

\[
\mathrm{Q} (1) = (r - 1)! \lim _ {n \rightarrow \infty} n ^ {1 - r} \alpha_ {n}.\tag{3}
\]

If \(F = \sum_{n \in \mathbf{Z}} a_n T^n\) and \(G = \sum_{n \in \mathbf{Z}} b_n T^n\) are two elements of \(Z((T))\) , we denote by “\(F \leqslant G\)” the relation “\(a_n \leqslant b_n\) for every \(n \in \mathbf{Z}\).” It is an order relation compatible with the ring structure of \(Z((T))\) (A, VI, p. 18, def. 1). We have \((1 - T)^{-1} \geqslant 1\) . If \(Q \in Z[T, T^{-1}]\) is \(\geqslant 0\) , then the integer \(Q(1)\) is positive.

Lemma 2. --- a) Let F be a nonzero element of \(\mathbf{Z}((\mathbf{T}))\) such that there exists \(r \in \mathbf{Z}\) , with \((1 - \mathrm{T})^r\mathrm{F} \in \mathbf{Z}[\mathrm{T}, \mathrm{T}^{-1}]\) ; then F is written uniquely in the form \(\mathrm{F} = (1 - \mathrm{T})^{-d}\) . Q, where \(\mathrm{Q} \in \mathbf{Z}[\mathrm{T}, \mathrm{T}^{-1}]\) , \(\mathrm{Q}(1) \neq 0\) and \(d \in \mathbf{Z}\) . If \(\mathrm{F} \geqslant 0\) , then we have \(\mathrm{Q}(1) > 0\) and \(d \geqslant 0\) .

\begin{enumerate}
\def\labelenumi{\alph{enumi})}
\setcounter{enumi}{1}
\tightlist
\item
  Let Q, R be in Z{[}T, T \(^{-1}\) {]}, and d, d' in Z with Q(1) \textgreater{} 0. If
\end{enumerate}

\[
(1 - \mathrm{T}) ^ {- d}. \mathrm{Q} \leqslant (1 - \mathrm{T}) ^ {- d ^ {\prime}}. \mathrm{R},
\]

then either \(d < d'\) , or \(d = d'\) and \(\mathbf{Q}(1) \leqslant \mathbf{R}(1)\) .

\begin{enumerate}
\def\labelenumi{\alph{enumi})}
\item
  One can write \(F = (1 - T)^{-r} T^{n} P(T)\) with r, \(n \in Z\) and \(P(T) \in Z[T]\) . By Euclidean division, one can write \(P(T) = (1 - T)^{p} R(T)\) with \(R(T) \in Z[T]\) and \(R(1) \neq 0\). Therefore \(F = (1 - T)^{-(r-p)} Q(T)\) , where \(Q(T) = T^{n} R(T) \in Z[T, T^{-1}]\) and \(Q(1) \neq 0\) . This proves the existence of d and Q. Moreover, if \((1 - T)^{r} Q(T) = (1 - T)^{s} R(T)\) with r \textgreater{} s and Q, R in \(Z[T, T^{-1}]\) , one has \(R(T) = (1 - T)^{r-s} Q(T)\) , hence \(R(1) = 0\) ; this proves uniqueness. Suppose that F is \(\geqslant 0\) ; if one had d \textless{} 0, then one would have \(F(1) = 0\) which is impossible since F is nonzero and all its coefficients are positive; we therefore have \(d \geqslant 0\) . If d = 0, then Q = F \(\geqslant 0\) , hence Q(1) is positive. If \(d \geqslant 1\) , then Q(1) is positive by formula (3). This proves a).
\item
  Suppose \(d \geqslant d'\) . Then \((1 - T)^{-d}((1 - T)^{d - d'} R - Q) \geqslant 0\) ; since \(S(T) = (1 - T)^{d - d'} R - Q\) belongs to \(Z[T, T^{-1}]\) , this implies \(S(1) \geqslant 0\) from what precedes. If \(d > d'\) , one has \(S(1) = -Q(1) < 0\) , whence a contradiction; if \(d = d'\) , one has \(S(1) = R(1) - Q(1)\) whence \(Q(1) \leqslant R(1)\) .
\end{enumerate}

